
\documentclass[12pt,a4paper]{article}

\usepackage[margin=1in]{geometry}
\usepackage{amsmath,amssymb,amsthm}
\usepackage{mathtools}
\usepackage[hidelinks]{hyperref}
\usepackage{booktabs}
\usepackage{array}
\usepackage{enumitem}
\usepackage{xcolor}
\usepackage{fontspec}

\defaultfontfeatures{Ligatures=TeX}
\IfFontExistsTF{Times New Roman}
  {\setmainfont{Times New Roman}}
  {\setmainfont{TeX Gyre Termes}}
\IfFontExistsTF{Menlo}
  {\setmonofont{Menlo}}
  {\setmonofont{Latin Modern Mono}}

\newcolumntype{P}[1]{>{\raggedright\arraybackslash}p{#1}}

\newtheorem{theorem}{Theorem}[section]
\newtheorem{proposition}[theorem]{Proposition}
\newtheorem{lemma}[theorem]{Lemma}
\newtheorem{corollary}[theorem]{Corollary}
\theoremstyle{definition}
\newtheorem{definition}[theorem]{Definition}
\newtheorem{assumption}[theorem]{Assumption}
\theoremstyle{remark}
\newtheorem{remark}[theorem]{Remark}

\newcommand{\tightlist}{%
  \setlength{\itemsep}{0pt}\setlength{\parskip}{0pt}}

\title{%
  \textbf{The Null Cluster as Structural Imprint}\\
  \textbf{Recurring Null-Boundary Convergence}\\
  \large A Programme Note after Q-RAIF A%
}

\author{
  Sidong Liu, PhD \\
  \small iBioStratix Ltd \\
  \small \texttt{sidongliu@hotmail.com}
}

\date{May 2026\\
\small Zenodo DOI: \href{https://doi.org/10.5281/zenodo.20283057}{10.5281/zenodo.20283057}}

\begin{document}
\emergencystretch=2em
\raggedbottom
\hbadness=5000
\vbadness=5000

\maketitle

\begin{abstract}
Several major structures in relativistic physics are organized at or near null boundaries: causal cones, event and Rindler horizons, horizon thermality, covariant entropy bounds, asymptotic symmetries, soft theorems, memory effects, and the propagation of massless fields. These facts are well known separately. The present note proposes a programme-level unification: these phenomena form a null cluster, a repeated convergence pattern in which causal reachability, thermal response, entropy localization, infrared memory, modular structure, and symmetry constraints co-localize at the null boundary of Lorentzian read-out.

The claim is not that any one of these phenomena is newly derived here. Nor is it that photon or graviton masslessness, BMS symmetry, Unruh thermality, or holographic entropy are new predictions. The contribution is structural: within the Q-RAIF A programme, where admissible finite read-out from algebraic data induces an effective Lorentzian geometry, the null boundary is read as the thinnest interface of the emergent geometry. It is the place where finite access, causal boundary, state restriction, entropy accounting, and asymptotic imprint most tightly meet.

This note therefore functions as a programme note rather than a proof. It identifies the null cluster, explains why its repeated appearance is not a trivial observation, places it within the Q-RAIF A and mass-energy recovery setting, and states failure modes that would weaken the proposed reading.

\medskip
\noindent\textbf{Keywords}: Q-RAIF A, null boundaries, Lorentzian read-out, horizon thermality, BMS symmetry, soft theorems, memory effect, modular flow, structural imprint
\end{abstract}

\clearpage
\tableofcontents
\clearpage
\section{Introduction}
\label{sec:introduction}

This note has a narrow purpose.

It does not introduce a new field equation, a new particle, or a new
quantum-gravity model. It does not claim to derive special relativity,
the Standard Model mass spectrum, black-hole entropy, soft theorems, or
the cosmological constant. It also does not replace ordinary general
relativity, quantum field theory, algebraic QFT, or scattering theory.

Its purpose is different:

\begin{quote}
to identify a repeated structural convergence around null boundaries and
to propose that this convergence is naturally read, within Q-RAIF A, as
a structural imprint of Lorentzian read-out.
\end{quote}

The central term is \textbf{null cluster}.

By the null cluster I mean the recurring co-localization of the
following features:

\begin{enumerate}[leftmargin=*]
\tightlist
\item
  causal reachability,
\item
  horizon or boundary structure,
\item
  thermal response under restricted access,
\item
  entropy localization on null or near-null sheets,
\item
  infrared memory and asymptotic symmetry,
\item
  modular flow with geometric meaning in special cases,
\item
  massless or near-massless propagation as a kinematic locking to the
  null cone.
\end{enumerate}

Not every feature appears in every physical situation. The claim is
weaker and more useful:

\begin{quote}
whenever a physical regime is organized by a null boundary, several of
these features tend to reappear together.
\end{quote}

This is not a theorem. It is a programme-level organizing principle.

\section{What Is Already Known}\label{what-is-already-known}

The individual ingredients of the null cluster are standard.

The causal cone of Lorentzian geometry is defined by the null
directions. Event horizons are null hypersurfaces. Rindler horizons
produce observer-dependent thermality. Black-hole entropy scales with
horizon area. Bousso's covariant entropy bound uses null light-sheets.
Asymptotically flat gravity has null infinity, BMS symmetries, soft
graviton theorems, and gravitational memory. In algebraic QFT, the
Bisognano-Wichmann theorem connects modular flow for wedge algebras with
Lorentz boosts. In high-energy or massless limits, propagation is locked
to the null cone.

None of that is new.

The useful question is different:

\begin{quote}
Why do these structures repeatedly meet at null boundaries?
\end{quote}

Mainstream literature usually studies partial clusters. The references
in the table are representative anchors, not exhaustive bibliographies.

\begin{center}
\begin{tabular}{P{0.21\linewidth}P{0.20\linewidth}P{0.23\linewidth}P{0.20\linewidth}}
\toprule
\textbf{Line of work} & \textbf{Representative anchor} & \textbf{Null phenomena emphasized} & \textbf{Usually outside its main focus} \\
\midrule
BMS / infrared triangle & Strominger 2018 & Asymptotic symmetry, soft theorems, memory & Horizon thermality, entropy bounds, modular flow \\
Thermodynamic gravity & Jacobson 1995 & Local horizons, heat, Einstein equation & BMS, soft theorems, memory \\
Covariant entropy bounds & Bousso 1999 & Null light-sheets, entropy limits & Modular theory, BMS, massless protection \\
Algebraic QFT / modular theory & Bisognano--Wichmann 1975; Haag 1996 & Wedge modular flow, Rindler thermality & BMS, covariant entropy bounds \\
Conformal / twistor approaches & Penrose 2010 & Null geometry, conformal structure & Horizon thermality, entropy localization \\
Gauge and diffeomorphism protection & Weinberg 1965; Wald 1984 & Massless propagation, constraint structure & Horizon thermality, entropy bounds, soft memory \\
\bottomrule
\end{tabular}
\end{center}

The present note proposes that these partial clusters are not merely
parallel accidents. They are different faces of a broader null cluster.

\section{Why the Observation Is Not
Trivial}\label{why-the-observation-is-not-trivial}

It is tempting to say: ``Of course these things involve light cones.
This is just relativity.''

That response misses the point.

The nontrivial point is not that light cones exist. The nontrivial point
is that multiple independent-looking structures concentrate at the same
kind of boundary:

\begin{itemize}
\tightlist
\item
  causal structure says what can influence what;
\item
  horizon thermality says restricted access to a wedge or exterior
  region produces a thermal state;
\item
  entropy bounds say information is controlled by area or null sheets,
  not by naive volume;
\item
  soft theorems and memory say radiation leaves persistent boundary data
  at null infinity;
\item
  modular theory says, in important cases, state-dependent algebraic
  flow can acquire geometric boost meaning;
\item
  massless propagation says protected gauge or diffeomorphism structures
  lock certain excitations to the null cone.
\end{itemize}

These are not the same statement in different notation.

They come from different parts of physics:

\begin{itemize}
\tightlist
\item
  Lorentzian geometry,
\item
  general relativity,
\item
  quantum field theory,
\item
  algebraic QFT,
\item
  black-hole thermodynamics,
\item
  scattering theory,
\item
  asymptotic symmetry,
\item
  gauge theory.
\end{itemize}

The priority value of this note is to draw these lines on one map and
ask whether their common null locus is a structural clue.

\section{The Null Cluster}\label{the-null-cluster}

\subsection{Core Cluster}\label{core-cluster}

The core cluster consists of phenomena that directly occur on, or are
defined by, null boundaries.

\begin{center}
\begin{tabular}{P{0.19\linewidth}P{0.25\linewidth}P{0.19\linewidth}P{0.22\linewidth}}
\toprule
\textbf{Node} & \textbf{Standard meaning} & \textbf{Null locus} & \textbf{Programme reading} \\
\midrule
Causal cone & Boundary between timelike and spacelike separation & $ds^2=0$ & Basic boundary of effective read-out \\
Event horizon & Boundary of causal access & Null hypersurface & Finite access boundary in geometry \\
Rindler horizon & Acceleration horizon & Null planes & Restricted wedge access yields thermality \\
Horizon thermality & Unruh / Hawking type thermal response & Horizon-local & State restriction at null boundary \\
Entropy bounds & Area or light-sheet control of entropy & Null sheets / horizons & Finite read-out closes on boundary data \\
BMS / soft / memory & Asymptotic symmetry and persistent radiation imprint & Null infinity & Null boundary stores scattering imprint \\
Modular geometry & Modular flow becomes boost in special QFT settings & Wedge / horizon structure & Algebraic flow can carry geometric boundary meaning \\
\bottomrule
\end{tabular}
\end{center}

These are the strongest candidates for the null cluster proper.

\subsection{Adjacent Node: Massless
Protection}\label{adjacent-node-massless-protection}

Massless protection should not be treated as identical to the core
cluster.

Photon masslessness is protected by gauge structure. Graviton
masslessness, in the usual low-energy picture, is tied to diffeomorphism
structure and massless spin-2 consistency. These are not the same
mechanisms as horizon thermality, entropy bounds, or BMS memory.

Still, they belong near the cluster.

The reason is \textbf{kinematic locking}:

\begin{quote}
gauge and diffeomorphism structures protect the absence of mass terms,
and the absence of mass locks propagation to the Lorentzian null cone.
\end{quote}

Thus massless protection is not a core null-boundary phenomenon in the
same sense as a horizon or null infinity. It is an adjacent node: a
symmetry-protected route by which field propagation becomes null.

This distinction matters. If massless protection is placed inside the
core cluster too quickly, the note overclaims. If it is excluded
completely, the relation between symmetry protection and null
propagation is lost.

The correct placement is second ring, not first ring.

\subsection{Inclusion and Exclusion
Criteria}\label{inclusion-and-exclusion-criteria}

The cluster is not meant to include every use of null geometry.

A phenomenon belongs to the \textbf{core cluster} only if it is defined
on, defined by, or operationally localized at a null boundary: an event
horizon, a Rindler horizon, a null light-sheet, null infinity, or a
wedge/null-plane structure whose algebraic meaning depends on that
boundary.

A phenomenon belongs to the \textbf{adjacent ring} if it is not itself a
null-boundary phenomenon, but is locked to the null cone by a symmetry
or kinematic mechanism. Massless protection is the main example: gauge
and diffeomorphism structures do not equal horizon thermality or BMS
memory, but they constrain propagation to the null cone.

A phenomenon is \textbf{outside the cluster} if it merely uses null
structure as a technical tool without making a null boundary carry
causal, thermal, entropic, asymptotic, modular, or imprint data. Twistor
methods, conformal techniques, causal-set discretizations, and generic
light-cone coordinates may become relevant in later work, but they are
not automatically part of the null cluster as defined here.

This criterion is meant to block cherry-picking. The cluster is not
``everything involving light cones''; it is the narrower convergence of
boundary, access, entropy, memory, modular, and symmetry-locking
structures around null loci.

\section{The Structural-Imprint
Claim}\label{the-structural-imprint-claim}

The programme-level claim can now be stated.

\begin{quote}
The null cluster is a structural imprint of Lorentzian read-out: it
marks the boundary where finite causal access, state restriction,
entropy accounting, asymptotic radiation data, and symmetry-protected
massless propagation most tightly meet.
\end{quote}

In Q-RAIF A, the effective Lorentzian metric is not taken as a brute
background. It is treated as a read-out structure induced under
admissibility constraints from the algebraic source pair. The
mass-energy companion paper then shows, conditionally, that once this
Lorentzian read-out and an admissible four-momentum bridge are in place,
the standard mass-energy relation is recovered inside the read-out.

The present note takes a further programme-level step:

if Lorentzian read-out is the arena, then its null boundary should be
where the most compressed structural traces appear.

This is why the same class of features keeps recurring:

\begin{itemize}
\tightlist
\item
  causal closure,
\item
  horizon restriction,
\item
  thermal response,
\item
  entropy localization,
\item
  soft boundary data,
\item
  memory,
\item
  modular-geometric flow,
\item
  massless propagation.
\end{itemize}

The null boundary is not merely the path of light. It is the boundary
grammar of Lorentzian manifestation.

\section{Pairwise Structural Links}\label{pairwise-structural-links}

The null cluster becomes more visible when the pairwise links are made
explicit.

\subsection{Causal Boundary and
Horizon}\label{causal-boundary-and-horizon}

A horizon is not a material wall. It is a causal boundary. Its
definition depends on which signals can reach which region, and this is
organized by null curves or null hypersurfaces.

Thus the causal cone and horizon structure are not separate topics. A
horizon is what the causal cone becomes when access is globally or
regionally restricted.

\subsection{Horizon and Thermality}\label{horizon-and-thermality}

The Unruh and Hawking effects show that horizon-related access
restrictions are tied to thermal response.

In a Rindler wedge, the restriction of the vacuum to the wedge is
thermal with respect to boost time. In black-hole settings, the horizon
is tied to Hawking temperature. The details differ, but the repeated
pattern is the same: restricted access at a horizon gives a thermal
reading.

Programme reading:

\begin{quote}
thermal response is not added on top of a horizon; it is one of the
signatures of finite access to a null-bounded region.
\end{quote}

\subsection{Horizon and Entropy}\label{horizon-and-entropy}

Black-hole entropy and covariant entropy bounds indicate that entropy is
controlled by boundary geometry, especially null or horizon-like
structure.

The important shift is from volume to boundary.

Programme reading:

\begin{quote}
when finite read-out closes a region, the entropy-relevant data are
organized at the boundary of access, not throughout a naive container
volume.
\end{quote}

The quantum focusing conjecture and related quantum Bousso-bound work
sharpen this line by replacing purely classical area focusing with
generalized entropy and quantum expansion along null congruences. This
does not prove the present note's imprint reading, but it reinforces the
point that entropy constraints and null deformation data are not
separate bookkeeping devices.

\subsection{Null Infinity and
Memory}\label{null-infinity-and-memory}

In scattering theory, radiation reaches null infinity. Soft theorems,
BMS symmetries, and memory effects show that null infinity is not empty
bookkeeping. It carries persistent imprint.

Programme reading:

\begin{quote}
null infinity is where scattering leaves boundary-readable memory.
\end{quote}

This is the asymptotic analogue of the horizon claim. Both cases treat
null boundaries as places where finite read-out becomes constrained by
persistent boundary data.

Recent work on BMS charge algebras, double-soft theorems, BMS fluxes,
and celestial holography makes this point more concrete: null infinity
is not merely a formal endpoint of radiation, but a boundary arena where
flux, symmetry, memory, and candidate celestial operator data are
organized.

\subsection{Modular Flow and Null
Boundaries}\label{modular-flow-and-null-boundaries}

The Bisognano-Wichmann theorem shows that, in relativistic QFT under
suitable assumptions, the modular flow of a wedge algebra is implemented
by Lorentz boosts.

This does not mean modular flow is always physical time, nor that
modular flow directly defines momentum.

The narrower lesson is enough:

\begin{quote}
in important cases, state-dependent algebraic flow knows about geometric
wedge and horizon structure.
\end{quote}

This is precisely the kind of link Q-RAIF A treats as structurally
significant, while still avoiding the overclaim that modular flow alone
derives all spacetime dynamics.

The limitation is important. The Bisognano-Wichmann result is a
statement about wedge algebras in relativistic QFT under specific
structural assumptions; it is not a generic theorem about arbitrary null
surfaces, arbitrary states, or curved spacetime. Interacting theories,
Haag duality issues, Type III factors, and the detailed role of Connes
spectrum belong to a deeper AQFT analysis that this note does not
attempt.

There is, however, targeted recent work closer to the present boundary
language: studies of modular operators for null-plane algebras in free
fields decompose the one-particle structure into lightlike fibres and
relate null cuts, relative entropy, and QNEC-type statements. This
supports the limited claim made here: modular data can have genuine
null-boundary content, but the technical domain of validity must be
stated explicitly.

\subsection{Massless Propagation and Null
Cone}\label{massless-propagation-and-null-cone}

A massless excitation propagates on the null cone. In gauge theory, the
absence of a photon mass is protected by gauge invariance. In gravity,
the masslessness of the low-energy spin-2 mode is tied to diffeomorphism
structure and consistency constraints.

Programme reading:

\begin{quote}
massless propagation is the field-theoretic way a symmetry-protected
degree of freedom is locked to the Lorentzian boundary.
\end{quote}

This is adjacent to, but not identical with, the horizon and entropy
parts of the cluster.

\section{Placement after Q-RAIF A}\label{placement-after-q-raif-a}

Q-RAIF A supplies the prior claim used here:

\begin{quote}
admissible finite read-out induces effective Lorentzian geometry under
the stated algebraic constraints.
\end{quote}

The mass-energy companion paper adds:

\begin{quote}
once admissible four-momentum is defined on the effective read-out
sector, the standard mass-energy relation is recovered as a conditional
theorem.
\end{quote}

This note adds a programme-level map:

\begin{quote}
once Lorentzian read-out is available, the null boundary becomes the
natural interface where causal, thermal, entropic, asymptotic, modular,
and massless-propagation features repeatedly concentrate.
\end{quote}

The three statements have different epistemic status.

\begin{center}
\begin{tabular}{P{0.44\linewidth}P{0.46\linewidth}}
\toprule
\textbf{Claim} & \textbf{Status} \\
\midrule
Q-RAIF A Lorentzian read-out & Prior theorem in the programme \\
Mass-energy recovery & Conditional theorem given Lorentzian read-out and momentum bridge \\
Null cluster as structural imprint & Programme-level unifying note \\
\bottomrule
\end{tabular}
\end{center}

This separation is essential. The present note should not be cited as
proving the null cluster. It proposes the null cluster as a structural
reading and a research map.

\section{What Is New}\label{what-is-new}

The note does not claim novelty for the following:

\begin{itemize}
\tightlist
\item
  light cones,
\item
  event horizons,
\item
  Hawking or Unruh thermality,
\item
  black-hole entropy,
\item
  covariant entropy bounds,
\item
  BMS symmetry,
\item
  soft theorems,
\item
  memory effects,
\item
  massless photon or graviton propagation,
\item
  modular theory in AQFT.
\end{itemize}

The novelty is the cross-domain placement.

The note says:

\begin{enumerate}[leftmargin=*]
\tightlist
\item
  these are not merely isolated null-related facts;
\item
  they form a repeated cluster around null boundaries;
\item
  Q-RAIF A gives a natural programme-level reason to expect such
  clustering, because null boundaries are the thinnest boundary of
  Lorentzian read-out;
\item
  the cluster suggests specific failure modes and future directions.
\end{enumerate}

This is a perspective contribution, not a new calculation.

\section{Failure Modes}\label{failure-modes}

The following would weaken or falsify the proposed reading.

\subsection{F1. Fundamental vacuum
dispersion}\label{f1.-fundamental-vacuum-dispersion}

If future observations establish a fundamental, non-environmental,
energy-dependent speed of massless propagation in vacuum, then the idea
of a single effective Lorentzian null cone would be weakened.

This would not automatically refute every algebraic-emergence programme,
but it would force Q-RAIF A to be modified toward multi-cone,
Finsler-like, Lorentz-violating, or energy-dependent read-out geometry.

The relevant target is not ordinary plasma dispersion, gravitational
Shapiro delay, source-intrinsic time lag, or effective propagation in
matter. The dangerous signal would be a vacuum law of the form

\[
v(E)=c\left[1+O\!\left(\left(\frac{E}{E_{\mathrm{QG}}}\right)^n\right)\right]
\]

or another non-environmental Lorentz-violating dispersion that survives
multi-messenger and multi-frequency controls. In such a case the note's
single-null-cone reading would have to be weakened or replaced.

\subsection{F2. Basic photon mass}\label{f2.-basic-photon-mass}

If the photon is found to have a fundamental nonzero rest mass, not
merely an effective mass in a medium or plasma, then the role of the
photon as the cleanest electromagnetic null pointer would fail.

The broader null cluster could survive, but its most familiar massless
representative would be downgraded.

\subsection{F3. Basic graviton mass}\label{f3.-basic-graviton-mass}

If gravitational waves are shown to arise from a fundamentally massive
spin-2 mode, rather than a massless geometric perturbation or an
effective modification, then the placement of gravitational waves as
direct null-boundary geometry would fail.

\subsection{F4. Mass-shell failure inside a valid Lorentzian
sector}\label{f4.-mass-shell-failure-inside-a-valid-lorentzian-sector}

If a system has:

\begin{itemize}
\tightlist
\item
  a valid effective Lorentzian geometry,
\item
  a well-defined four-momentum,
\item
  and ordinary energy-momentum read-out,
\end{itemize}

but violates the standard mass-shell relation in a way that cannot be
attributed to medium effects, interactions, curved-background
corrections, or bad variable choice, then the mass-energy companion
bridge would be undermined.

\subsection{F5. No deep relation among horizon thermality, entropy,
memory, and modular
structure}\label{f5.-no-deep-relation-among-horizon-thermality-entropy-memory-and-modular-structure}

If future work showed that horizon thermality, entropy localization,
infrared memory, and modular geometry are only unrelated special-case
coincidences, then the null cluster would lose its unifying force.

This is the most likely soft failure mode. It would not be a single
experimental refutation, but a gradual failure of the programme to
generate useful connections.

\section{Open Directions}\label{open-directions}

\subsection{Null Boundary as Imprint
Interface}\label{null-boundary-as-imprint-interface}

The most publishable next technical direction is to connect BMS
symmetry, soft theorems, memory effects, and imprint language.

The guiding idea:

\begin{quote}
null infinity is not only where radiation escapes; it is where
scattering leaves boundary-readable imprint.
\end{quote}

This can potentially be written as a technical essay with standard
references and without extending Q-RAIF machinery.

\subsection{Cosmological Constant as Global Read-out
Curvature}\label{cosmological-constant-as-global-read-out-curvature}

A second direction is the cosmological constant.

The cautious claim:

\begin{quote}
the cosmological constant problem may be partly a category problem if
one treats global read-out curvature as ordinary local vacuum energy
density.
\end{quote}

This is not a replacement for Lambda-CDM. It is not a solution to the
numerical problem. It is a programme-level reframing.

\subsection{Black-Hole Information as Read-out Sector
Reallocation}\label{black-hole-information-as-read-out-sector-reallocation}

The black-hole information problem can be cautiously read as a mismatch
between sector-limited external read-out and global algebraic
accounting.

Modern island and replica-wormhole results already suggest that the
accessible description of entropy depends on how the gravitational and
radiation sectors are assigned.

The programme-level reading:

\begin{quote}
the question is not simply where information goes, but which read-out
sector can access which imprint.
\end{quote}

\subsection{Neutrinos near the Null
Boundary}\label{neutrinos-near-the-null-boundary}

Neutrino masses are extremely small compared with other known fermion
masses. This makes them natural candidates for a ``near-null''
structural reading.

This should remain a weak direction. Seesaw mechanisms, Majorana versus
Dirac mass, sterile neutrino constraints, and mass ordering are concrete
physics questions. The programme should not replace them with metaphor.

The only safe statement is:

\begin{quote}
neutrinos may be an example where a tiny timelike thickness sits very
close to the null class.
\end{quote}

\subsection{Dark Sector as Limited Observable
Access}\label{dark-sector-as-limited-observable-access}

The dark sector is the riskiest direction.

It is tempting to say that dark matter is not a particle but a read-out
effect. That is too strong.

The publishable version is narrower:

\begin{quote}
some dark-sector phenomena may be usefully interpreted as effects of
limited observable access, where gravitational read-out is broader than
electromagnetic read-out.
\end{quote}

This remains an open conjecture. It requires quantitative modelling
before it can become a prediction.

\section{Limits}\label{limits}

This note has strict limits.

\begin{enumerate}[leftmargin=*]
\tightlist
\item
  It does not derive the null cluster.
\item
  It does not derive BMS symmetry, soft theorems, memory effects,
  horizon entropy, or modular flow.
\item
  It does not claim that all null-related phenomena have the same
  physical mechanism.
\item
  It does not treat gauge invariance and diffeomorphism invariance as
  the same mechanism.
\item
  It does not claim that dark matter is not a particle.
\item
  It does not solve the cosmological constant problem.
\item
  It does not solve black-hole information.
\item
  It does not introduce religious or metaphysical terminology into the
  physics note.
\end{enumerate}

The correct label is:

\begin{quote}
structural programme note.
\end{quote}

\section{Minimal Claim}\label{minimal-claim}

The minimal claim is this:

Null boundaries repeatedly gather causal limits, horizon restriction,
thermal response, entropy localization, asymptotic memory, modular
geometry, and symmetry-protected massless propagation.

Q-RAIF A gives this recurrence a programme-level placement: the null
boundary is the thinnest interface of Lorentzian read-out, where finite
access and boundary data meet.

This is not a new equation; it is a map.

% ============================================================
\begin{thebibliography}{99}

\bibitem{Einstein1905}
A.~Einstein, \emph{On the Electrodynamics of Moving Bodies} (1905).

\bibitem{Wald1984}
R.~M.~Wald, \emph{General Relativity}, University of Chicago Press (1984).

\bibitem{Hawking1975}
S.~W.~Hawking, \emph{Particle creation by black holes}, Communications in Mathematical Physics \textbf{43}, 199--220 (1975).

\bibitem{Bekenstein1973}
J.~D.~Bekenstein, \emph{Black holes and entropy}, Physical Review D \textbf{7}, 2333--2346 (1973).

\bibitem{Unruh1976}
W.~G.~Unruh, \emph{Notes on black-hole evaporation}, Physical Review D \textbf{14}, 870--892 (1976).

\bibitem{Bousso1999}
R.~Bousso, \emph{A covariant entropy conjecture}, Journal of High Energy Physics \textbf{07}, 004 (1999).

\bibitem{BoussoFisherLeichenauerWall2016}
R.~Bousso, Z.~Fisher, S.~Leichenauer, and A.~C.~Wall, \emph{Quantum focusing conjecture}, Physical Review D \textbf{93}, 064044 (2016).

\bibitem{BoussoFisherKoellerLeichenauerWall2016}
R.~Bousso, Z.~Fisher, J.~Koeller, S.~Leichenauer, and A.~C.~Wall, \emph{Proof of the quantum null energy condition}, Physical Review D \textbf{93}, 024017 (2016).

\bibitem{Jacobson1995}
T.~Jacobson, \emph{Thermodynamics of spacetime: The Einstein equation of state}, Physical Review Letters \textbf{75}, 1260--1263 (1995).

\bibitem{BondiMetznerSachs1962}
H.~Bondi, M.~G.~J.~van der Burg, and A.~W.~K.~Metzner, \emph{Gravitational waves in general relativity VII}, Proceedings of the Royal Society A \textbf{269}, 21--52 (1962).

\bibitem{Sachs1962}
R.~K.~Sachs, \emph{Gravitational waves in general relativity VIII}, Proceedings of the Royal Society A \textbf{270}, 103--126 (1962).

\bibitem{Weinberg1965}
S.~Weinberg, \emph{Infrared photons and gravitons}, Physical Review \textbf{140}, B516--B524 (1965).

\bibitem{Strominger2018}
A.~Strominger, \emph{Lectures on the Infrared Structure of Gravity and Gauge Theory}, Princeton University Press (2018).

\bibitem{CampigliaLaddha2021}
M.~Campiglia and A.~Laddha, \emph{BMS algebra, double soft theorems, and all that}, arXiv:2106.14717.

\bibitem{DonnayFiorucciHerfrayRuzziconi2021}
L.~Donnay, A.~Fiorucci, Y.~Herfray, and R.~Ruzziconi, \emph{BMS flux algebra in celestial holography}, arXiv:2108.11969.

\bibitem{Pasterski2021}
S.~Pasterski, \emph{Lectures on celestial amplitudes}, arXiv:2108.04801.

\bibitem{BisognanoWichmann1975}
J.~J.~Bisognano and E.~H.~Wichmann, \emph{On the duality condition for a Hermitian scalar field}, Journal of Mathematical Physics \textbf{16}, 985--1007 (1975).

\bibitem{Haag1996}
R.~Haag, \emph{Local Quantum Physics: Fields, Particles, Algebras}, Springer-Verlag (1996).

\bibitem{Takesaki2003}
M.~Takesaki, \emph{Theory of Operator Algebras II}, Springer (2003).

\bibitem{MorinelliTanimotoWegener2021}
V.~Morinelli, Y.~Tanimoto, and B.~Wegener, \emph{Modular operator for null plane algebras in free fields}, arXiv:2107.00039.

\bibitem{Penrose2010}
R.~Penrose, \emph{Cycles of Time}, Bodley Head (2010).

\bibitem{Liu2026MassEnergy}
S.~Liu, \emph{Recovering Mass--Energy Equivalence within Algebraic Lorentzian Read-out: A Conditional Theorem in Q-RAIF A}, Zenodo (2026), DOI: 10.5281/zenodo.20280951.

\end{thebibliography}

\end{document}
